\scanpage{163}
\margin{\handwritten{\uline{Nouveau fascicule}}}

\begin{center}
\textsc{Chapitre Premier}

LE LANGAGE DES SCHÉMAS

\textbf{Sommaire}
\end{center}

\noindent\S~1. Schémas affines.

\noindent\S~2. Préschémas et morphismes de préschémas.

\noindent\S~3. Produits de préschémas.

\noindent\S~4. Sous-préschémas et morphismes d’immersion.

\noindent\S~5. Préschémas réduits ; condition de séparation.

\noindent\S~6. Conditions de finitude.

\noindent\S~7. Applications rationnelles.

\noindent\S~8. Les schémas de Chevalley.

\noindent\S~9. Compléments sur les faisceaux quasi cohérents.

\noindent\S~10. Schémas formels.

Les \S\S~1 à 8 ne font guère que développer un langage, celui qui sera utilisé dans toute la suite. Notons cependant que, conformément à l’esprit général de ce Traité, les \S\S~7 et 8 seront moins utilisés que les autres, et de façon moins essentielle ; on n’a d’ailleurs parlé des schémas de Chevalley que pour faire le lien avec le langage de Chevalley [1] et Nagata [9]. Le \S~9 donne des définitions et résultats sur les faisceaux quasi cohérents, dont certains ne se bornent plus à une traduction en langage « géométrique » de notions connues d’Algèbre commutative, mais sont déjà de nature globale ; ils seront indispensables, dès les chapitres suivants, dans l’étude globale des morphismes. Enfin, le \S~10 introduit une généralisation de la notion de schéma, qui nous servira d’intermédiaire au chapitre III pour formuler et démontrer de façon commode les résultats fondamentaux de l’étude cohomologique des morphismes propres ; par ailleurs, signalons que la notion de schéma formel semble indispensable pour exprimer certains faits de la « théorie des modules » (problèmes de classification des variétés algébriques). Les résultats du \S~10 ne seront pas utilisés avant le \S~3 du chapitre III et il est recommandé d’en omettre la lecture jusque-là.
% archive sheet number: 162
